@@M0@@ (VI 5.5), one
may suppose that @@M1@@ is an invertible sheaf. We then have: @@M2@@; hence @@M3@@.

The @@M4@@, @@M5@@ vanish for @@M6@@. It therefore suffices to show that @@M7@@. Now @@M8@@ is the class in @@M9@@ of @@M10@@, and @@M11@@ that of @@M12@@. The equality:
@@M13@@
gives:
@@M14@@
the result follows.

\sgaStableTarget{sga6.e08.thm.3-6}\textbf{Theorem 3.6. (Riemann--Roch).} \uline{Let} @@M15@@ \uline{be a projective complete-intersection morphism} (1.1), \uline{whose virtual relative dimension} @@M16@@ \uline{is constant} (1.9). \uline{Suppose that} @@M17@@ \uline{is quasi-compact}, \uline{and has an ample sheaf}. \uline{With the notation of} 3.3, \uline{we then have for every} @@M18@@ \uline{the relation}:
@@M19@@
@@M20@@ \uline{denoting the dual element of the class of the relative cotangent complex} @@M21@@ (2.5 \uline{and} 2.3).

In other words, the diagram
@@M22@@
is commutative.
